\documentclass[10pt,a4paper]{amsart}
\usepackage[margin=19mm]{geometry}
\usepackage{amsmath,amssymb,mathtools}
\usepackage[scr=rsfs,cal=euler]{mathalfa}
\usepackage{xfrac}
\usepackage[unicode,hidelinks,bookmarks=false]{hyperref}
\newcommand{\ie}{i.e.\ }
\newcommand{\cf}{cf.\ }
\newcommand{\resp}{respectively\ }
\newcommand{\Subsection}{\subsection}
\providecommand{\nobreakdash}{\nobreak\hbox{-}}
\setlength{\emergencystretch}{2em}
\multlinegap0pt
\usepackage{bm}
\usepackage{bbm}
\usepackage{dsfont}
\usepackage{xspace}
\usepackage{cancel}
\usepackage{mathtools}
\usepackage{amsthm}
\makeatletter
\@ifundefined{thm}{\newtheorem{thm}{Thm}}{}
\@ifundefined{lemma}{\newtheorem{lemma}[thm]{Lemma}}{}
\makeatother
\DeclareMathOperator{\nil}{\mathcal N}
\title[Profinite groups with many elements with large nilpotentizer]{Profinite groups with many elements with large nilpotentizer and generalizations}
\author{Martino Garonzi, Andrea Lucchini, Nowras Otmen}
\date{Source: arXiv:2505.16589v1 (CC BY 4.0)}
\begin{document}
\maketitle

\noindent\emph{Source and licence.} Profinite groups with many elements with large nilpotentizer and generalizations. Version arXiv:2505.16589v1 (CC BY 4.0), distributed under the Creative Commons Attribution 4.0 International licence, as recorded in the arXiv licence metadata for this exact version and shown on the abstract page https://arxiv.org/abs/2505.16589v1. Full rights notice: licenses/arxiv-2505.16589-CC-BY-4.0.txt.\par\smallskip

\noindent\textit{Editorial scope.} Counted unit: the Lemma whose source label is \texttt{nilabcf} in the article (source0.tex lines 514--519, source label \texttt{nilabcf}), delivered together with the article's own complete original proof of it (source0.tex lines 521--538, one \texttt{proof} environment of 18 lines). No mathematics is written, repaired or re-derived: statement and proof are the article's own text line for line. Required context retained uncounted: none. Every definition, display and label of those lines is retained unchanged. Omitted from this package and not redistributed: 1--513, 520--520, 539--779. Nothing inside a retained excerpt is omitted or altered: every retained body is delivered exactly as the article's lines are written, so no omission receipt of any kind accompanies this package. Citations: the retained excerpts carry no citation command at all, so no bibliography entry of the article is transcribed and no reference is redistributed. Reference adaptation: none was needed and none was made; every reference inside the delivered unit resolves inside it, so no unresolved reference or citation is produced. Printed slips kept literal: no printed word, symbol, hypothesis or number of the counted statement or of its proof was altered. Layout and class: the article's own class is \texttt{amsart} with packages including babel, inputenc, amsmath, amssymb, epic, graphicx, mathrsfs, enumerate, xy, xcolor, comment, enumitem, hyperref, cleveref; the wrapper is the lane's pinned \texttt{amsart} editorial wrapper at 10pt on A4, which restates the theorem-like environments the retained text needs and adds the article's own macro definitions that the retained text needs (\texttt{\textbackslash nil}), with extra wrapper packages (bm, bbm, dsfont, xspace, cancel, mathtools). The delivered numbering is the unit's own. Assets: no retained span contains a figure environment, a graphics-inclusion command, a PDF-page inclusion, a table or a TikZ picture; no image file is copied into this package, the package declares no asset dependency and no image file is redistributed with it. Licence: version arXiv:2505.16589v1 of this article is distributed under the Creative Commons Attribution 4.0 International licence, as recorded in the arXiv licence metadata for this exact version and shown on the abstract page https://arxiv.org/abs/2505.16589v1. A full rights notice is delivered at licenses/arxiv-2505.16589-CC-BY-4.0.txt. Characters: every retained excerpt is pure ASCII, so the delivered engine, XeLaTeX with its default Latin Modern text font, reports no missing character.
\begin{lemma}\label{nilabcf}
	Let $G$ be a finite group, let $p$ be a prime and suppose that $N$ is an abelian minimal normal subgroup of $G$ whose order is a $p$-power. Let $x\in G$ and write
	$x=x_1x_2$ where $|x_1|$ is a $p$-power, $|x_2|$ is coprime with $p$ and $[x_1,x_2]=1.$
	Then
	$$\frac{|\nil_G(x)|}{|G|}\leq \frac{|\nil_{G/N}(xN)|}{|G/N|}\frac {|C_N(x_2)|}{|N|}.$$
\end{lemma}

\begin{proof}
   Given $y\in \nil_G(x)$ we want to estimate the size of $\Delta:=\nil_G(x)\cap Ny.$ We write $x=x_1x_2$ and $y=y_1y_2$, where $|x_1|$, $|y_1|$ are $p$-powers, $|x_2|, |y_2|$ are coprime with $p$ and $[x_1,x_2]=[y_1,y_2]=1.$ Since $\langle x,y
	\rangle$ is nilpotent, $[x_2,y_1]=[x_1,y_2]=1.$  Now assume $z\in \Delta.$ Then there exist $n_1, n_2\in N$ such that $z=y_1n_1y_2n_2,$ where $|y_1n_1|$ is a $p$-power, $|y_2n_2|$ is coprime with $p$ and $[y_1n_1,y_2n_2]=1.$
	First notice that $H = \langle x_2,y_2 \rangle$ and $K = \langle x_2, y_2n_2\rangle$ are Hall $p'$-subgroups of the solvable group $\langle x_2,y_2 \rangle N$, so by Hall's theorem there exists $n \in N$ such that $H^n = K$. Since $N \cap K = \{1\}$, if $h \in H$ then $|hN \cap K| \leq 1$, moreover $h^n \in hN$, so ${x_2}^n = x_2$ and ${y_2}^n = y_2n_2$. In other words, $y_2n_2 = {y_2}^{n}$ for some
	$n\in C_N(x_2).$ Thus the number of possible choices for
	$y_2n_2$ is at most $|C_N(x_2)|/|C_N(x_2)\cap C_N(y_2)|.$ Moreover
	$\langle x_2, y_1n_1\rangle$ nilpotent implies that $n_1\in C_N(x_2).$ Finally, using the commutator equalities $[ab,c]=[a,c]^b [b,c]$ and $[c,ab]=[c,b][c,a]^b$, since $N$ is abelian we have 
    \begin{align*}
        1 & = [y_1n_1,y_2n_2] = [y_1,n_2]^{n_1} [y_1,y_2]^{n_2n_1} [n_1,n_2] [n_1,y_2]^{n_2} = [y_1,n_2][n_1,y_2],
    \end{align*}
    thus $[y_1,n_2]=[y_2,n_1].$ Once $n_2$ is fixed, we count how many $n_1$ there are in $C_N(x_2)$ with $[y_1,n_2]=[y_2,n_1].$ It is possible that there is no $n_1$ satisfying this equation. In any case assume that 
	$[y_2,n_1]=[y_2,\tilde n_1]$, where $n_1$ and $\tilde n_1$ are two different elements of $C_N(x_2).$ Then $n_1\tilde n_1^{-1}\in C_N(y_2)\cap C_N(x_2)$, so, given $n_2$ there are at most $|C_N(x_2)\cap C_N(y_2)|$ possible choices for $n_1$. We conclude
	that the number of possibilities for the pair $(n_1,n_2)$
	is at most
	$$\frac{|C_N(x_2)|}{|C_N(x_2)\cap C_N(y_2)|}\cdot |C_N(x_2)\cap C_N(y_2)|=|C_N(x_2)|.
	$$
	Hence $|\Delta|\leq |C_N(x_2)|$. This implies the result.
\end{proof}

\end{document}
